\documentclass[11pt,a4paper]{article}

% ── Codificación, fuentes y lenguaje ───────────────────────────────────────────
\usepackage[utf8]{inputenc}
\usepackage[T1]{fontenc}
\usepackage[default,scale=0.95]{sourcesanspro}
\renewcommand{\familydefault}{\sfdefault}
\usepackage[spanish,es-noshorthands,es-tabla]{babel}

% ── Geometría y espaciado ─────────────────────────────────────────────────────
\usepackage[top=2.2cm, bottom=2.5cm, left=2.2cm, right=2.2cm, headheight=15pt]{geometry}
\usepackage{setspace}
\setstretch{1.18}
\usepackage{parskip}

% ── Matemáticas y unidades ────────────────────────────────────────────────────
\usepackage{amsmath}
\usepackage{amssymb}
\usepackage{siunitx}

% ── Circuitos ─────────────────────────────────────────────────────────────────
\usepackage[european, straightvoltages]{circuitikz}

% ── Tablas ────────────────────────────────────────────────────────────────────
\usepackage{booktabs}
\usepackage{tabularx}
\usepackage{array}
\usepackage{longtable}
\usepackage{needspace}
\usepackage{colortbl}
\usepackage{enumitem}

% ── Colores, cajas e iconos ───────────────────────────────────────────────────
\usepackage[dvipsnames]{xcolor}
\usepackage{tcolorbox}
\tcbuselibrary{skins, breakable}
\usepackage{fontawesome5}
\usepackage{graphicx}

% ── Listados de código (dark-mode infográfico) ────────────────────────────────
%   Estilo base: fondo oscuro con números de línea. Los bloques lstlisting se
%   envuelven automáticamente en una caja tcolorbox redondeada.
\usepackage{listings}

% ── Secciones con barra de color (infografía) ────────────────────────────────
\usepackage{titlesec}
\titleformat{\section}
  {\normalfont\LARGE\bfseries\color{azulNoche}}
  {\colorbox{cyanNeon}{\color{fondoOscuro}\thesection}}{0.7em}{}
  [\vspace{2pt}\color{cyanNeon}\rule{\linewidth}{1.8pt}]
\titlespacing*{\section}{0pt}{24pt}{10pt}
\titleformat{\subsection}
  {\normalfont\large\bfseries\color{azulNoche}}
  {\textcolor{cyanNeon}{\thesubsection}}{0.7em}{}
  [\vspace{1pt}\color{grisLinea}\rule{\linewidth}{0.6pt}]
\titlespacing*{\subsection}{0pt}{14pt}{6pt}
\titleformat{\subsubsection}
  {\normalfont\normalsize\bfseries\color{azulNoche}}
  {\textcolor{cyanNeon}{\thesubsubsection}}{0.7em}{}
\titlespacing*{\subsubsection}{0pt}{10pt}{4pt}

% ── Cabeceras y pies de página ────────────────────────────────────────────────
\usepackage{fancyhdr}
\usepackage{lastpage}
\pagestyle{fancy}
\fancyhf{}
\renewcommand{\headrulewidth}{0pt}
\renewcommand{\footrulewidth}{0pt}
\fancyfoot[C]{%
  \begin{minipage}{\linewidth}
    \vspace{2pt}
    \color{cyanNeon}\rule{\linewidth}{1.2pt}\\[2pt]
    \centering\color{grisTexto}\footnotesize
    Página \thepage\ de \pageref{LastPage}
  \end{minipage}%
}

% ── Hiperenlaces ──────────────────────────────────────────────────────────────
\usepackage[colorlinks=true, linkcolor=azulNoche, urlcolor=cyanNeon]{hyperref}
\sloppy
\emergencystretch 3em

% ── Paleta de colores — tecnológico dark-mode ──────────────────────────────────
\definecolor{fondoOscuro}{HTML}{0D1B2A}%             Dark navy — fondo banda título
\definecolor{verdeTurquesa}{HTML}{004D40}%           Verde turquesa — bandas de marca
\definecolor{azulNoche}{HTML}{1B3A6B}%               Midnight blue — secciones, marcos
\definecolor{azulMedio}{HTML}{2A5298}%               Azul medio — variante de acento
\definecolor{cyanNeon}{HTML}{00C8E0}%                Cyan eléctrico — acentos primarios
\definecolor{verdeSignal}{HTML}{00B887}%             Verde señal — fortalezas, éxito
\definecolor{naranjaVivo}{HTML}{FF7043}%             Naranja vivo — mejoras, advertencias
\definecolor{rojoAlerta}{HTML}{E53935}%              Rojo alerta — errores
\definecolor{grisPapel}{HTML}{F4F6FA}%               Fondo tarjetas claras
\definecolor{grisLinea}{HTML}{C8D0E0}%               Bordes sutiles
\definecolor{grisTexto}{HTML}{6B7A99}%               Texto secundario
\definecolor{amarilloNota}{HTML}{FFB300}%            Notas / advertencias doradas
\definecolor{codigoFondo}{HTML}{1E2D3D}%             Fondo del bloque de código
\definecolor{codigoComentario}{HTML}{7A8FA6}%        Comentarios
\definecolor{codigoCadena}{HTML}{FF9D5C}%            Cadenas / strings
\definecolor{codigoKeyword}{HTML}{00C8E0}%           Palabras clave
\definecolor{codigoNumero}{HTML}{00E5A0}%            Números

% ── Banda de título: portada infográfica ──────────────────────────────────────
%   Diseño en 2 capas: banda de color (por defecto fondo oscuro) + franja de
%   acento cyan abajo. Uso: \bandaTitulo[color]{título}{subtítulo}.
%   El icono se puede cambiar con \renewcommand{\iconoBanda}{...} (FontAwesome).
\newcommand{\iconoBanda}{microchip}
\newcommand{\bandaTitulo}[3][fondoOscuro]{%
  \begin{tcolorbox}[
    enhanced,
    colback=#1, colframe=#1,
    boxrule=0pt, arc=8pt, outer arc=8pt,
    width=\linewidth,
    top=18pt, bottom=6pt, left=14pt, right=14pt,
    borderline south={3.5pt}{0pt}{cyanNeon},
  ]
    \begin{minipage}[c]{0.07\linewidth}
      \centering
      \textcolor{cyanNeon}{\fontsize{30}{30}\selectfont\faIcon{\iconoBanda}}
    \end{minipage}%
    \hspace{8pt}%
    \begin{minipage}[c]{0.88\linewidth}
      {\fontsize{20}{24}\selectfont\bfseries\color{white}#2}\par\vspace{5pt}%
      \textcolor{cyanNeon!80!white}{\small\faIcon{angle-right}~#3}%
    \end{minipage}
    \vspace{4pt}
  \end{tcolorbox}%
  \vspace{8pt}%
}

% ── Tarjetas de datos (infografía) ────────────────────────────────────────────
\tcbset{
  tarjetaDato/.style={
    enhanced, breakable,
    arc=6pt, outer arc=6pt,
    colback=grisPapel, colframe=azulNoche,
    boxrule=1pt,
    fonttitle=\bfseries\small, coltitle=white,
    attach boxed title to top left={yshift=-3mm, xshift=6mm},
    boxed title style={
      colback=azulNoche, colframe=azulNoche,
      arc=4pt, boxrule=0pt,
    },
    drop fuzzy shadow=azulNoche!30!white,
    top=8pt, bottom=8pt, left=10pt, right=10pt,
  },
  tarjetaIcono/.style={
    enhanced, breakable,
    arc=6pt, outer arc=6pt,
    colback=white, colframe=grisLinea, boxrule=0.8pt,
    fonttitle=\bfseries\small, coltitle=azulNoche,
    borderline west={3pt}{0pt}{cyanNeon},
    drop fuzzy shadow=azulNoche!25!white,
    top=6pt, bottom=6pt, left=10pt, right=10pt,
  },
  cajaTitulo/.style={
    enhanced, breakable,
    arc=6pt, outer arc=6pt,
    colback=azulNoche!10!grisPapel, colframe=azulNoche,
    boxrule=1pt,
    fonttitle=\bfseries, coltitle=white,
    attach boxed title to top left={yshift=-3mm, xshift=6mm},
    boxed title style={
      colback=azulNoche, colframe=azulNoche,
      arc=4pt, boxrule=0pt,
    },
    drop fuzzy shadow=azulNoche!25!white,
    top=8pt, bottom=8pt, left=10pt, right=10pt,
  },
  cajaContenido/.style={
    enhanced, breakable,
    arc=5pt, outer arc=5pt,
    colback=grisPapel, colframe=grisLinea,
    boxrule=0.8pt,
    fonttitle=\bfseries\small, coltitle=azulNoche,
    top=6pt, bottom=6pt, left=10pt, right=10pt,
  },
  cajaFortaleza/.style={
    enhanced, breakable,
    arc=6pt, outer arc=6pt,
    colback=verdeSignal!8!white, colframe=verdeSignal,
    boxrule=1pt,
    borderline west={4pt}{0pt}{verdeSignal},
    fonttitle=\bfseries\small, coltitle=verdeSignal!70!black,
    drop fuzzy shadow=verdeSignal!20!white,
    top=6pt, bottom=6pt, left=10pt, right=10pt,
  },
  cajaMejora/.style={
    enhanced, breakable,
    arc=6pt, outer arc=6pt,
    colback=naranjaVivo!8!white, colframe=naranjaVivo,
    boxrule=1pt,
    borderline west={4pt}{0pt}{naranjaVivo},
    fonttitle=\bfseries\small, coltitle=naranjaVivo!80!black,
    drop fuzzy shadow=naranjaVivo!20!white,
    top=6pt, bottom=6pt, left=10pt, right=10pt,
  },
  cajaRecomendacion/.style={
    enhanced, breakable,
    arc=6pt, outer arc=6pt,
    colback=cyanNeon!6!white, colframe=cyanNeon!60!azulNoche,
    boxrule=1pt,
    borderline west={4pt}{0pt}{cyanNeon},
    fonttitle=\bfseries\small, coltitle=azulNoche,
    drop fuzzy shadow=azulNoche!20!white,
    top=6pt, bottom=6pt, left=10pt, right=10pt,
  },
}

% ── Ícono + texto (encabezados de sección y elementos) ────────────────────────
\newcommand{\iconotexto}[2]{\textcolor{cyanNeon}{\faIcon{#1}}~#2}

% ── Listados de código: estilo dark + auto-envoltura ─────────────────────────
\lstdefinestyle{estiloCodigo}{
    backgroundcolor=\color{codigoFondo},
    basicstyle=\ttfamily\small\color{white},
    commentstyle=\color{codigoComentario}\itshape,
    stringstyle=\color{codigoCadena},
    keywordstyle=\color{codigoKeyword}\bfseries,
    numberstyle=\tiny\color{codigoComentario},
    numbers=left,
    numbersep=10pt,
    showstringspaces=false,
    breaklines=true,
    breakatwhitespace=false,
    tabsize=4,
    frame=none,
    captionpos=b,
    aboveskip=6pt,
    belowskip=4pt,
    xleftmargin=14pt,
    framexleftmargin=14pt,
    literate=
      {á}{{\'a}}1 {é}{{\'e}}1 {í}{{\'i}}1 {ó}{{\'o}}1 {ú}{{\'u}}1
      {Á}{{\'A}}1 {É}{{\'E}}1 {Í}{{\'I}}1 {Ó}{{\'O}}1 {Ú}{{\'U}}1
      {ñ}{{\~n}}1 {Ñ}{{\~N}}1 {ü}{{\"u}}1 {Ü}{{\"U}}1
      {¿}{{?`}}1 {¡}{{!`}}1
}
\lstdefinestyle{estiloPython}{
    style=estiloCodigo,
    language=Python,
}
\lstset{style=estiloCodigo}
\BeforeBeginEnvironment{lstlisting}{%
  \begin{tcolorbox}[
    enhanced, breakable,
    arc=6pt, outer arc=6pt,
    colback=codigoFondo, colframe=azulNoche,
    boxrule=1pt,
    drop fuzzy shadow=azulNoche!25!white,
    top=2pt, bottom=2pt, left=0pt, right=0pt,
  ]%
}
\AfterEndEnvironment{lstlisting}{\end{tcolorbox}}

% ── Cajas de contenido temáticas ──────────────────────────────────────────────
\newtcolorbox{cajaRecuerda}{
    enhanced, breakable,
    arc=6pt, outer arc=6pt,
    colback=azulNoche!10!grisPapel,
    colframe=azulNoche, boxrule=1pt,
    borderline west={4pt}{0pt}{cyanNeon},
    fonttitle=\bfseries\small\color{white},
    title={\faIcon{info-circle}~Recuerda},
    attach boxed title to top left={yshift=-3mm, xshift=6mm},
    boxed title style={colback=azulNoche, arc=4pt, boxrule=0pt},
    drop fuzzy shadow=azulNoche!25!white,
    left=10pt, right=10pt, top=8pt, bottom=8pt,
}

\newtcolorbox{cajaConcepto}{
    enhanced, breakable,
    arc=6pt, outer arc=6pt,
    colback=verdeSignal!8!white,
    colframe=verdeSignal, boxrule=1pt,
    borderline west={4pt}{0pt}{verdeSignal},
    fonttitle=\bfseries\small\color{white},
    title={\faIcon{lightbulb}~Concepto clave},
    attach boxed title to top left={yshift=-3mm, xshift=6mm},
    boxed title style={colback=verdeSignal!80!black, arc=4pt, boxrule=0pt},
    drop fuzzy shadow=verdeSignal!20!white,
    left=10pt, right=10pt, top=8pt, bottom=8pt,
}

\newtcolorbox{cajaEjemplo}{
    enhanced, breakable,
    arc=6pt, outer arc=6pt,
    colback=cyanNeon!5!white,
    colframe=cyanNeon!70!azulNoche, boxrule=1pt,
    borderline west={4pt}{0pt}{cyanNeon},
    fonttitle=\bfseries\small\color{fondoOscuro},
    title={\faIcon{code}~Ejemplo},
    attach boxed title to top left={yshift=-3mm, xshift=6mm},
    boxed title style={colback=cyanNeon, arc=4pt, boxrule=0pt},
    drop fuzzy shadow=azulNoche!20!white,
    left=10pt, right=10pt, top=8pt, bottom=8pt,
}

\newtcolorbox{cajaEjercicio}{
    enhanced, breakable,
    arc=6pt, outer arc=6pt,
    colback=azulMedio!7!white,
    colframe=azulMedio, boxrule=1pt,
    borderline west={4pt}{0pt}{azulMedio},
    fonttitle=\bfseries\small\color{white},
    title={\faIcon{pencil-alt}~Ejercicio},
    attach boxed title to top left={yshift=-3mm, xshift=6mm},
    boxed title style={colback=azulMedio, arc=4pt, boxrule=0pt},
    drop fuzzy shadow=azulNoche!20!white,
    left=10pt, right=10pt, top=8pt, bottom=8pt,
}

\newtcolorbox{cajaReto}{
    enhanced, breakable,
    arc=6pt, outer arc=6pt,
    colback=naranjaVivo!6!white,
    colframe=naranjaVivo, boxrule=1pt,
    borderline west={4pt}{0pt}{naranjaVivo},
    fonttitle=\bfseries\small\color{white},
    title={\faIcon{fire}~Reto},
    attach boxed title to top left={yshift=-3mm, xshift=6mm},
    boxed title style={colback=naranjaVivo, arc=4pt, boxrule=0pt},
    drop fuzzy shadow=naranjaVivo!20!white,
    left=10pt, right=10pt, top=8pt, bottom=8pt,
}

% ── Indicadores de dificultad ─────────────────────────────────────────────────
\newcommand{\facil}{\textcolor{verdeSignal}{\faIcon{circle}\,\textbf{Fácil}}}
\newcommand{\intermedio}{\textcolor{naranjaVivo}{\faIcon{circle}\,\textbf{Intermedio}}}
\newcommand{\dificil}{\textcolor{rojoAlerta}{\faIcon{circle}\,\textbf{Difícil}}}

% ─────────────────────────────────────────────────────────────────────────────
\renewcommand{\iconoBanda}{flask}
\fancyhead[L]{\footnotesize\color{grisTexto}\faIcon{flask}~Laboratorio II de Física}
\fancyhead[R]{\footnotesize\color{grisTexto}Práctica 04 --- CIRCUITOS ELÉCTRICOS}
\hypersetup{pdftitle={Práctica 04 --- CIRCUITOS ELÉCTRICOS},pdfauthor={Prof. César Rodríguez \textbar{} Universidad de Oriente, Núcleo de Sucre, Departamento de Física \textbar{} Cumaná, abril de 2011}}

\begin{document}
\bandaTitulo[fondoOscuro]{Práctica 04 --- CIRCUITOS ELÉCTRICOS}{Prácticas según el programa de la Universidad de Oriente}

\begin{center}\small\color{grisTexto}Prof. César Rodríguez \textbar{} Universidad de Oriente, Núcleo de Sucre, Departamento de Física \textbar{} Cumaná, abril de 2011\end{center}

\begin{tcolorbox}[cajaContenido]\textbf{Fuente:} enunciado del PDF original (Universidad de Oriente, Núcleo de Sucre, Departamento de Física, 2011).
\end{tcolorbox}

\section{\iconotexto{bullseye}{Objetivos}}

\begin{itemize}[leftmargin=1.4em,itemsep=2pt,topsep=3pt]
  \item Medir corrientes y voltajes en circuitos conectados en serie, paralelo y mixto.
  \item Aplicar las leyes de Kirchhoff en redes eléctricas.
\end{itemize}

\section{\iconotexto{book-open}{Fundamento teórico}}

Dos o más resistencias conectadas de modo que la misma carga fluya a través de cada una de ellas se dice que están conectadas en serie; por ella debe circular la misma corriente $i$.

\begin{center}
\begin{circuitikz}
  \draw (0,0) to[V, l=$E$, invert] (0,3);
  \draw (0,3) -- (1,3);
  \draw (1,3) to[R, l=$R_1$, v=$V_1$, i=$I_1$] (3.5,3);
  \draw (3.5,3) to[R, l=$R_2$, v=$V_2$, i=$I_2$] (6,3);
  \draw (6,3) -- (6,0) -- (0,0);
\end{circuitikz}
\par\smallskip{\small\itshape Asociación de dos resistencias en serie}
\end{center}

**Figura 1. Dos resistencias conectadas en serie por las cuales circula** **una corriente de intensidad I.**

Para este tipo de conexión deben cumplirse las siguientes ecuaciones:

\begin{center}\small
\begin{tabular}{lll}
\toprule
Corriente: & $I = I_1 = I_2$ & (1.a) \\
\midrule
Voltaje: & $V_{ab} = V_1 + V_2$ & (1.b) \\
Resistencia: & $R_{eq} = R_1 + R_2$ & (1.c) \\
\bottomrule
\end{tabular}
\end{center}

Dos resistencias unidas como en la figura 2 se dice que están conectadas en paralelo. Sea $i$ la corriente que va desde "a" hasta "b". En el punto "a", la corriente se divide en dos, $i_1$ que circula por $R_1$ e $i_2$ por $R_2$.

\begin{center}
\begin{circuitikz}
  \draw (0,0) to[V, l=$E$, invert] (0,3);
  \draw (0,3) -- (1.5,3);
  \draw (1.5,3) to[R, l=$R_1$, v=$V_1$, i=$I_1$] (1.5,0);
  \draw (1.5,0) -- (0,0);
  \draw (1.5,3) -- (4,3);
  \draw (4,3) to[R, l=$R_2$, v=$V_2$, i=$I_2$] (4,0);
  \draw (4,0) -- (1.5,0);
\end{circuitikz}
\par\smallskip{\small\itshape Asociación de dos resistencias en paralelo}
\end{center}

**Figura 2. Dos resistencias conectadas en paralelo, por las cuales circulan** **las corrientes $I_1$ e $I_2$.**

Para este tipo de conexión deben cumplirse las siguientes ecuaciones:

\begin{center}\small
\begin{tabular}{lll}
\toprule
Corriente: & $I = I_1 + I_2$ & (2.a) \\
\midrule
Voltaje: & $V_{ab} = V_1 = V_2$ & (2.b) \\
Resistencia: & $R_{eq} = \frac{R_1 R_2}{R_1 + R_2}$ & (2.c) \\
\bottomrule
\end{tabular}
\end{center}

A veces no es posible reducir muchas de las redes importantes de electrónica a simples combinaciones serie-paralelo, por lo que deben utilizarse métodos analíticos más avanzados. A tal fin, resultan muy útiles las conocidas leyes de Kirchhoff. Estas leyes pueden ser enunciadas de la siguiente manera:

**Primera ley:** En cualquier nodo de un circuito, la suma de las corrientes que llegan a él es igual a la suma de las corrientes que salen del mismo, es decir, la suma algebraica de las corrientes en un nodo es igual a cero. \[ \sum I_i = 0 \] (3) **Segunda ley:** En una malla, la suma algebraica de los productos IR y de las f.e.m que se encuentran en un recorrido es igual a cero. \[ \sum IR + \sum \epsilon_i = 0 \] (4)

Para aplicar esta Ley es necesario definir:

\begin{itemize}[leftmargin=1.4em,itemsep=2pt,topsep=3pt]
  \item **Nodo:** punto donde se unen tres o más conductores.
  \item **Malla:** Cualquier recorrido seguido de un circuito cerrado.
  \item **Signos algebraicos e intensidades:** se acostumbra tomar como positivas
\end{itemize}

las corrientes que llegan a un nodo y negativas las que salen de él.

\begin{itemize}[leftmargin=1.4em,itemsep=2pt,topsep=3pt]
  \item **Signos algebraicos de voltajes:** Al recorrer una malla en cualquier
\end{itemize}

dirección se encontrarán resistencias y baterías. En las resistencias las diferencias de potencial son positivas si se atraviesan de los puntos de mayor a menor potencial y negativas si se atraviesan de menor a mayor potencial.

\begin{itemize}[leftmargin=1.4em,itemsep=2pt,topsep=3pt]
  \item La f.e.m. batería o pila: es positiva si se recorre del borne (-) al borne (+) y
\end{itemize}

negativa si se recorre de (+) a (-).

Estos arreglos son convencionales, pero una vez que se haga la selección de una convención de signos determinada ésta debe seguirse a través de todo el problema.

**Reglas para la aplicación correcta de las leyes de Kirchhoff:**

\begin{itemize}[leftmargin=1.4em,itemsep=2pt,topsep=3pt]
  \item Dibuje un diagrama de la red, claro y de tamaño suficiente. Use letras y
\end{itemize}

símbolos para cada corriente, resistencia o voltaje desconocido.

\begin{itemize}[leftmargin=1.4em,itemsep=2pt,topsep=3pt]
  \item Indique el sentido de cada corriente con una flecha. Si el sentido no lo conoce,
\end{itemize}

suponga momentáneamente un sentido en cualquier dirección. Si supone errado el sentido, la corriente le aparecerá con signo (-) al resolver las ecuaciones.

\begin{itemize}[leftmargin=1.4em,itemsep=2pt,topsep=3pt]
  \item Coloque un signo (+) al terminal positivo de cada batería y al extremo de cada
\end{itemize}

resistencia por donde usted supone que la corriente entra.

\begin{itemize}[leftmargin=1.4em,itemsep=2pt,topsep=3pt]
  \item Si tiene "n" nodos entonces habrá (n-1) ecuaciones independientes.
  \item Si "m" es el número de mallas, entonces usted tendrá (m-1) ecuaciones
\end{itemize}

independientes.

\section{\iconotexto{box-open}{Materiales y equipos}}

\begin{itemize}[leftmargin=1.4em,itemsep=2pt,topsep=3pt]
  \item 1 Tablero para redes.
  \item 6 pilas tipo D.
  \item 3 Resistencias de carbón.
  \item 1 multimedidor Keithley modelo 179 o 175.
\end{itemize}

\section{\iconotexto{list-ol}{Procedimiento}}

Se usa una red en la cual hay conectadas varias baterías y cajas de resistencias. Se determinan las intensidades de corriente utilizando valores conocidos de las resistencias y de los voltajes de las baterías.

\subsection*{RESISTENCIAS EN SERIE}

\begin{itemize}[leftmargin=1.4em,itemsep=2pt,topsep=3pt]
  \item Mida con el voltímetro la f.e.m. y con el Ohmmétro los valores de R$_{1}$ y R$_{2}$. Organice estos datos en una tabla.
  \item Arme el circuito mostrado en la figura 3.
  \item Mida las corrientes (I, I$_{1}$ e I$_{2}$) con el amperímetro.
  \item Mida la diferencia de potencial entre los extremos de cada resistencia.
\end{itemize}

\begin{center}
\begin{circuitikz}
  \draw (0,0) to[V, l=$E$, invert] (0,3.5);
  \draw (0,3.5) to[switch, l=Int.] (1.5,3.5);
  \draw (1.5,3.5) to[ammeter, l=$I$] (3,3.5);
  \draw (3,3.5) to[R, l=$R_1$] (5,3.5);
  \draw (5,3.5) to[R, l=$R_2$] (7,3.5);
  \draw (7,3.5) -- (7,0) -- (0,0);
  \draw (3,3.5) -- (3,2) to[voltmeter, l=$V_1$] (5,2) -- (5,3.5);
  \draw (5,3.5) -- (5,1) to[voltmeter, l=$V_2$] (7,1) -- (7,3.5);
\end{circuitikz}
\par\smallskip{\small\itshape Medida de tensiones y corriente en el circuito serie}
\end{center}

Figura 3. Circuito experimental.

\subsection*{RESISTENCIAS EN PARALELO}

\begin{itemize}[leftmargin=1.4em,itemsep=2pt,topsep=3pt]
  \item Mida con el voltímetro la f.e.m. $\epsilon$ y con el Ohmmétro los valores de R$_{1}$ y R$_{2}$. Organice estos datos en una tabla.
  \item Arme el circuito mostrado en la figura 4.
  \item Mida la corriente total (I) la que circula por cada resistencia (I$_{1}$ e I$_{2}$).
  \item Mida la diferencia de potencial entre los extremos de cada resistencia.
\end{itemize}

\begin{center}
\begin{circuitikz}
  \draw (0,0) to[V, l=$E$, invert] (0,4);
  \draw (0,4) to[switch, l=Int.] (1.5,4);
  \draw (1.5,4) to[ammeter, l=$I$] (3,4);
  \draw (3,4) -- (6,4);
  \draw (3,4) to[ammeter, l=$I_1$] (3,2.5) to[R, l=$R_1$] (3,0);
  \draw (6,4) to[ammeter, l=$I_2$] (6,2.5) to[R, l=$R_2$] (6,0);
  \draw (0,0) -- (6,0);
  \draw (3,2.5) -- (4.2,2.5);
  \draw (1.2,3.2) to[voltmeter, l=$V_1$] (1.2,0.8);
  \draw (1.2,3.2) -- (1.2,4);
  \draw (1.2,0.8) -- (1.2,0);
  \draw (7.5,4) -- (7.5,3.2) to[voltmeter, l=$V_2$] (7.5,0.8) -- (7.5,0);
  \draw (6,4) -- (7.5,4);
  \draw (6,0) -- (7.5,0);
\end{circuitikz}
\par\smallskip{\small\itshape Medida de tensiones y corrientes en el circuito paralelo}
\end{center}

Figura 4. Circuito experimental.

\subsection*{RED ELÉCTRICA}

\begin{itemize}[leftmargin=1.4em,itemsep=2pt,topsep=3pt]
  \item Disponga las resistencias y pilas en el tablero tal como se muestra en la figura 5.
  \item Mida con su error los valores de $\epsilon_{1}$, $\epsilon_{2}$, $\epsilon_{3}$, R$_{1}$, R$_{2}$ y R$_{3}$.
  \item Arme el correspondiente circuito.
  \item Mida la corriente de cada rama (i$_{1}$, i$_{2}$, e i$_{3}$); para esto, rompa una rama del circuito e intercale el amperímetro (en serie). Haga revisar las conexiones por el profesor.
  \item Mida para cada malla, la diferencia de potencial entre los extremos de cada resistencia
\end{itemize}

\begin{center}
\begin{circuitikz}
  \draw (0,0) to[V, l=$E_1$, invert] (0,2);
  \draw (0,2) to[R, l=$R_1$] (0,4) -- (0,4);
  \draw (0,4) to[ammeter, l=$I_1$] (2,4);
  \draw (2,4) -- (4,4);
  \draw (4,4) to[ammeter, l=$I_3$] (6,4);
  \draw (6,4) to[R, l_=$R_3$] (6,2);
  \draw (6,2) to[V, l_=$E_3$] (6,0);
  \draw (0,0) -- (6,0);
  \draw (2,4) to[ammeter, l=$I_2$] (2,2.6) to[R, l=$R_2$] (2,1.4) to[V, l=$E_2$, invert] (2,0);
\end{circuitikz}
\par\smallskip{\small\itshape Red con tres ramas y tres fuentes para las leyes de Kirchhoff}
\end{center}

Fig. 5. Circuito experimental de una red eléctrica

\section{\iconotexto{calculator}{Actividades}}

\subsection*{RESISTENCIAS EN SERIE}

\begin{itemize}[leftmargin=1.4em,itemsep=2pt,topsep=3pt]
  \item Calcule la corriente que circula por el circuito y compárela con el valor medido.
  \item Verifique si se cumple la ecuación [1.b].
\end{itemize}

\subsection*{RESISTENCIAS EN PARALELOS}

\begin{itemize}[leftmargin=1.4em,itemsep=2pt,topsep=3pt]
  \item Calcule la corriente total que circula por el circuito y las corrientes i$_{1}$ e i$_{2}$. Compárelas con el valor medido.
  \item Verifique si se cumple la ecuación (2.-a).
  \item Verifique si se cumple la ecuación (2-b).
\end{itemize}

\subsection*{RED ELÉCTRICA}

\begin{itemize}[leftmargin=1.4em,itemsep=2pt,topsep=3pt]
  \item Aplique las leyes de Kirchhoff. Resuelva las ecuaciones en las incógnitas i$_{1}$, i$_{2}$, e i$_{3}$. (Debe traerlo resuelto antes de entrar al laboratorio).
  \item Con los valores medidos de $\epsilon_{1}$, $\epsilon_{2}$, $\epsilon_{3}$, R$_{1}$, R$_{2}$ y R$_{3}$. Calcule los valores de i$_{1}$, i$_{2}$, e i$_{3}$
  \item Calcule el porcentaje de error entre los valores experimentales obtenidos para las corrientes y los calculados por las leyes de Kirchhoff.
  \item Verifique si se cumple la ecuación (3).
  \item Verifique si se cumple la ecuación (4)
\end{itemize}
\end{document}
